\section{Les lleis dels gasos}

L'aire és un \textbf{gas}, i els gasos es comporten d'una manera molt característica: si els \textbf{comprimim} s'escalfen i fan més pressió; si els \textbf{escalfem} es volen expandir. Aquestes regles són les \textbf{lleis dels gasos}, i expliquen per què funciona un compressor, per què s'escalfa una bomba de bicicleta i per què els dipòsits tenen una vàlvula de seguretat.

\begin{idea}
L'estat d'una quantitat de gas ($n$ constant, no en perdem ni n'afegim) es descriu amb tres magnituds:
\begin{itemize}
 \item la \textbf{pressió} $p$ (sempre \textbf{absoluta}, en bar o Pa),
 \item el \textbf{volum} $V$ (en L o m$^3$),
 \item la \textbf{temperatura} $T$ (sempre \textbf{en kelvin}: $T = t + 273$).
\end{itemize}
Si en canviem una, les altres reaccionen. Les lleis ens diuen \emph{com}.
\end{idea}

\subsection*{1. Llei de Boyle--Mariotte (temperatura constant)}

Una xeringa tapada pel dit: si empenys l'èmbol, el volum d'aire disminueix i la pressió augmenta, i viceversa.

\begin{center}
\begin{minipage}{0.52\linewidth}
\begin{idea}
Si la \textbf{temperatura és constant}, la pressió i el volum són \textbf{inversament proporcionals}:
\[ p_1 \cdot V_1 = p_2 \cdot V_2 \]
Si el volum es fa 2 vegades més petit, la pressió es fa 2 vegades més gran.
\end{idea}
\end{minipage}\hfill
\begin{minipage}{0.4\linewidth}\centering
\begin{tikzpicture}[x=0.9cm,y=0.9cm]
 \draw[-{Stealth}] (0,0)--(4.4,0) node[right]{\small $V$};
 \draw[-{Stealth}] (0,0)--(0,3.6) node[above]{\small $p$};
 \draw[blau,thick,domain=0.9:4.2,samples=40] plot (\x,{3.6/\x});
 \fill[vermell] (1,3.6) circle (2pt) node[above right,font=\scriptsize]{$p_1,V_1$};
 \fill[vermell] (3,1.2) circle (2pt) node[above right,font=\scriptsize]{$p_2,V_2$};
\end{tikzpicture}
\end{minipage}
\end{center}

\begin{exemple}{el dipòsit d'un compressor}
Un compressor agafa \SI{200}{\liter} d'aire de l'atmosfera (\SI{1}{\bar} absolut) i els comprimeix en un dipòsit de \SI{25}{\liter}, sense que la temperatura canviï. Quina és la pressió final?
\begin{itemize}
 \item \textbf{Dades:} $p_1 = \SI{1}{\bar}$, $V_1 = \SI{200}{\liter}$, $V_2 = \SI{25}{\liter}$. Busquem $p_2$.
 \item \textbf{Fórmula:} $p_1 V_1 = p_2 V_2 \ \Rightarrow\ p_2 = \dfrac{p_1 V_1}{V_2}$
 \item \textbf{Substitució:} $p_2 = \dfrac{1 \cdot 200}{25} = \SI{8}{\bar}$ \textbf{absoluts}, és a dir, \SI{7}{\bar} al manòmetre.
\end{itemize}
\end{exemple}

\subsection*{2. Llei de Gay-Lussac (volum constant)}

Un dipòsit rígid tancat, al sol, es posa calent. Com no pot canviar de volum, l'única que puja és la pressió.

\begin{idea}
Si el \textbf{volum és constant}, la pressió i la temperatura (en kelvin) són \textbf{directament proporcionals}:
\[ \frac{p_1}{T_1} = \frac{p_2}{T_2} \]
Si la temperatura absoluta es doblega, la pressió es doblega.
\end{idea}

\begin{exemple}{dipòsit escalfat}
Un dipòsit conté aire a \SI{6}{\bar} absoluts i \SI{20}{\degreeCelsius}. S'escalfa fins a \SI{60}{\degreeCelsius}. Quina pressió hi ha ara?
\begin{itemize}
 \item \textbf{Dades:} $p_1 = \SI{6}{\bar}$, $T_1 = 20 + 273 = \SI{293}{\kelvin}$, $T_2 = 60 + 273 = \SI{333}{\kelvin}$.
 \item \textbf{Fórmula:} $p_2 = p_1 \cdot \dfrac{T_2}{T_1}$
 \item \textbf{Substitució:} $p_2 = 6 \cdot \dfrac{333}{293} = \SI{6,82}{\bar}$
\end{itemize}
Per això els dipòsits tenen una \textbf{vàlvula de seguretat}: si s'escalfen massa, la pressió podria superar el límit.
\end{exemple}

\subsection*{3. Llei de Charles (pressió constant)}

Un globus deixat al sol s'infla; ficat a la nevera es desinfla. Si la pressió no canvia, el volum del gas varia amb la temperatura.

\begin{idea}
Si la \textbf{pressió és constant}, el volum i la temperatura (en kelvin) són \textbf{directament proporcionals}:
\[ \frac{V_1}{T_1} = \frac{V_2}{T_2} \]
\end{idea}

\begin{exemple}{gas que s'expandeix}
\SI{2}{\liter} de gas a \SI{27}{\degreeCelsius} s'escalfen fins a \SI{127}{\degreeCelsius} a pressió constant. Quin volum ocupen?
\begin{itemize}
 \item \textbf{Dades:} $V_1 = \SI{2}{\liter}$, $T_1 = 27 + 273 = \SI{300}{\kelvin}$, $T_2 = 127 + 273 = \SI{400}{\kelvin}$.
 \item \textbf{Fórmula:} $V_2 = V_1 \cdot \dfrac{T_2}{T_1}$
 \item \textbf{Substitució:} $V_2 = 2 \cdot \dfrac{400}{300} = \SI{2,67}{\liter}$
\end{itemize}
\end{exemple}

\subsection*{La llei general (les tres alhora)}

Quan canvien les tres magnituds, les tres lleis es combinen:
\[ \boxed{\ \frac{p_1 \cdot V_1}{T_1} = \frac{p_2 \cdot V_2}{T_2}\ } \]
Si una magnitud no canvia, desapareix de l'equació i obtens una de les tres lleis anteriors.

\begin{center}
\renewcommand{\arraystretch}{1.5}
\begin{tabular}{@{}lccl@{}}
\toprule
\textbf{Llei} & \textbf{Constant} & \textbf{Fórmula} & \textbf{Què passa}\\
\midrule
Boyle--Mariotte & $T$ & $p_1V_1 = p_2V_2$ & $\uparrow p \Rightarrow \downarrow V$\\
Gay-Lussac & $V$ & $p_1/T_1 = p_2/T_2$ & $\uparrow T \Rightarrow \uparrow p$\\
Charles & $p$ & $V_1/T_1 = V_2/T_2$ & $\uparrow T \Rightarrow \uparrow V$\\
\bottomrule
\end{tabular}
\end{center}

\begin{exemple}{llei general en un compressor}
Un compressor aspira \SI{100}{\liter} d'aire a \SI{1}{\bar} i \SI{20}{\degreeCelsius}, i els comprimeix fins a \SI{10}{\liter}. En comprimir-se, l'aire s'escalfa fins a \SI{80}{\degreeCelsius}. Quina és la pressió final?
\begin{itemize}
 \item \textbf{Dades:} $p_1 = \SI{1}{\bar}$, $V_1 = \SI{100}{\liter}$, $T_1 = \SI{293}{\kelvin}$, $V_2 = \SI{10}{\liter}$, $T_2 = \SI{353}{\kelvin}$.
 \item \textbf{Fórmula:} $p_2 = \dfrac{p_1 V_1 T_2}{T_1 V_2}$
 \item \textbf{Substitució:} $p_2 = \dfrac{1 \cdot 100 \cdot 353}{293 \cdot 10} = \SI{12,05}{\bar}$ absoluts.
\end{itemize}
Si no s'escalfés, la pressió seria \SI{10}{\bar} (Boyle). L'escalfament la fa pujar encara més.
\end{exemple}

\begin{metode}{resoldre un problema de gasos}
\begin{itemize}
 \item \textbf{\ding{172}} Identifica \textbf{què es manté constant} ($T$, $V$ o $p$) i escull la llei.
 \item \textbf{\ding{173}} Passa la \textbf{pressió a absoluta} i la \textbf{temperatura a kelvin}.
 \item \textbf{\ding{174}} Escriu la fórmula, aïlla la incògnita i substitueix.
 \item \textbf{\ding{175}} Comprova que el resultat té sentit (si comprimeixes, la pressió ha de pujar).
\end{itemize}
\end{metode}

\begin{atencio}
Dos errors molt freqüents: fer servir la pressió del manòmetre (relativa) en comptes de l'\textbf{absoluta}, i fer servir graus Celsius en comptes de \textbf{kelvin}. Amb \SI{20}{\degreeCelsius} i \SI{40}{\degreeCelsius} sembla que la temperatura es dobla, però en kelvin només puja de \SI{293}{\kelvin} a \SI{313}{\kelvin}.
\end{atencio}

\begin{guiada}{una xeringa tapada}
Una xeringa conté \SI{60}{\milli\liter} d'aire a \SI{1}{\bar} (absolut). Tapem la sortida i empenyem l'èmbol fins que el volum és de \SI{20}{\milli\liter}, sense que la temperatura canviï. Quina és la pressió final?
\pas{1}{Què es manté constant? Quina llei?}{\vspace{0.8cm}}
\pas{2}{Fórmula:}{\vspace{0.8cm}}
\pas{3}{Substitució:}{\vspace{0.8cm}}
\pas{4}{Resultat (en bar absolut):}{\vspace{0.8cm}}
\resposta{0pt}{$T$ constant: llei de Boyle--Mariotte. $p_2 = p_1V_1/V_2 = 1\cdot60/20 = \SI{3}{\bar}$ absoluts.}
\end{guiada}

\subsection*{Exercicis}

\begin{exercici}{1}
Un gas ocupa \SI{5}{\liter} a \SI{1}{\bar}. Es comprimeix fins a \SI{1}{\liter} a temperatura constant. Quina pressió té?
\resposta{2.4cm}{$p_2 = 1\cdot5/1 = \SI{5}{\bar}$.}
\end{exercici}

\begin{exercici}{1}
Un dipòsit rígid té aire a \SI{2}{\bar} i \SI{300}{\kelvin}. S'escalfa fins a \SI{450}{\kelvin}. Quina pressió hi ha?
\resposta{2.4cm}{$p_2 = 2\cdot450/300 = \SI{3}{\bar}$.}
\end{exercici}

\begin{exercici}{1}
\SI{3}{\liter} de gas a \SI{250}{\kelvin} s'escalfen fins a \SI{500}{\kelvin} a pressió constant. Quin volum ocupen?
\resposta{2.4cm}{$V_2 = 3\cdot500/250 = \SI{6}{\liter}$.}
\end{exercici}

\begin{exercici}{2}
Un dipòsit de \SI{50}{\liter} conté aire a \SI{8}{\bar} absoluts. Quants litres d'aire a \SI{1}{\bar} (a l'atmosfera) hi ha dins el dipòsit?
\resposta{2.6cm}{Boyle: $V_2 = p_1V_1/p_2 = 8\cdot50/1 = \SI{400}{\liter}$.}
\end{exercici}

\begin{exercici}{2}
Un dipòsit rígid té aire a \SI{6}{\bar} absoluts i \SI{15}{\degreeCelsius}. Al sol arriba a \SI{45}{\degreeCelsius}. Quina pressió absoluta té?
\resposta{3cm}{$T_1 = \SI{288}{\kelvin}$, $T_2 = \SI{318}{\kelvin}$. $p_2 = 6\cdot318/288 = \SI{6,63}{\bar}$.}
\end{exercici}

\begin{exercici}{2}
El manòmetre d'un dipòsit rígid marca \SI{5}{\bar} a \SI{20}{\degreeCelsius}. S'escalfa fins a \SI{50}{\degreeCelsius}. Què marcarà ara el manòmetre?
\resposta{3.6cm}{$p_1 = 5 + 1 = \SI{6}{\bar}$ absoluts; $T_1 = \SI{293}{\kelvin}$, $T_2 = \SI{323}{\kelvin}$. $p_2 = 6\cdot323/293 = \SI{6,61}{\bar}$ absoluts $\Rightarrow$ el manòmetre marca $6{,}61 - 1 = \SI{5,61}{\bar}$.}
\end{exercici}

\begin{exercici}{2}
Un globus de \SI{1,5}{\liter} a \SI{20}{\degreeCelsius} es posa al congelador, a \SI{-30}{\degreeCelsius}. Quin volum té (suposa que la pressió no canvia)?
\resposta{3cm}{$T_1 = \SI{293}{\kelvin}$, $T_2 = \SI{243}{\kelvin}$. $V_2 = 1{,}5\cdot243/293 = \SI{1,24}{\liter}$.}
\end{exercici}

\begin{exercici}{3}
\SI{20}{\liter} d'aire a \SI{1}{\bar} i \SI{27}{\degreeCelsius} es comprimeixen en un cilindre fins a \SI{4}{\liter} i \SI{7}{\bar} absoluts. A quina temperatura (en \si{\degreeCelsius}) arriba l'aire?
\resposta{3.6cm}{$T_1 = \SI{300}{\kelvin}$. $T_2 = \dfrac{p_2V_2T_1}{p_1V_1} = \dfrac{7\cdot4\cdot300}{1\cdot20} = \SI{420}{\kelvin} = \SI{147}{\degreeCelsius}$.}
\end{exercici}

\begin{exercici}{3}
Quina llei dels gasos (Boyle, Gay-Lussac o Charles) explica cadascuna d'aquestes situacions?\\
a) La bomba de bicicleta s'escalfa quan infles una roda. \ b) Un dipòsit d'aire comprimit explota si es deixa dins d'un cotxe al sol. \ c) Una xeringa tapada és més difícil de prémer com més l'has premuda.
\resposta{3cm}{a) Combinada (augmenta la pressió, disminueix el volum i puja la temperatura); b) Gay-Lussac (volum constant, puja $T$ i puja $p$); c) Boyle--Mariotte (disminueix el volum i puja la pressió).}
\end{exercici}

\begin{comprova}
Si comprimeixes un gas a temperatura constant, la seva pressió\ldots\ \ $\square$ augmenta \quad $\square$ disminueix \quad $\square$ no canvia
\resposta{0pt}{Augmenta (Boyle--Mariotte).}
\end{comprova}
